% Cells visited (blue) by the kernels of Section 4.5 on an 8 x 8 array.
% x = coordinate 0 (ctr_0) to the right, y = coordinate 1 (ctr_1) upwards.
\begin{tikzpicture}[x=0.34cm,y=0.34cm,font=\scriptsize]
  \newcommand{\ispanel}[4]{% x-shift, filter (TeX condition on \x,\y), title, code
    \begin{scope}[shift={(#1,0)}]
      \foreach \x in {0,...,7} \foreach \y in {0,...,7} {
        \pgfmathparse{#2 ? 1 : 0}
        \ifnum\pgfmathresult=1 \fill[psblue!50] (\x,\y) rectangle ++(1,1); \fi
      }
      \draw[gridline] (0,0) grid (8,8);
      \draw[domainline,line width=0.5pt] (0,0) rectangle (8,8);
      \node[align=center,font=\footnotesize] at (4,9.2) {#3};
      \node[align=center,font=\ttfamily\tiny] at (4,-1.5) {#4};
    \end{scope}}
  \ispanel{0}{(\x>=1) && (\x<=6) && (\y>=1) && (\y<=6)}
    {(a) ghost layers}{ghost\_layers=1\\(AUTO for Jacobi)}
  \ispanel{10.5}{(\x>=1) && (\x<=4)}
    {(b) per side}{ghost\_layers=\\{}[(1, 3), 0]}
  \ispanel{21}{(mod(\x,2)==0) && (\y>=3) && (\y<=4)}
    {(c) slice}{iteration\_slice=\\make\_slice[::2, 3:-3]}
  \ispanel{31.5}{(\y>=\x)}
    {(d) triangle}{iteration\_slice=\\make\_slice[:, ctr\_0:]}
  \draw[axisarrow] (0,-3.4) -- (3,-3.4) node[right] {$x$ (\texttt{ctr\_0})};
  \draw[axisarrow] (-0.6,0) -- (-0.6,3) node[above] {$y$};
\end{tikzpicture}
